\documentclass[onsided]{book}
\usepackage[backend=biber,natbib=true,style=authoryear]{biblatex}
\addbibresource{/home/nguyen/1_NQBH/math/bib.bib}
\usepackage[utf8]{inputenc}
\usepackage{float}
\usepackage{graphicx}
\usepackage[colorlinks=true,linkcolor=blue,urlcolor=red,citecolor=magenta]{hyperref}
\usepackage{amsmath,amssymb,amsthm}
\allowdisplaybreaks
\numberwithin{equation}{section}
\newtheorem{assumption}{Assumption}[section]
\newtheorem{lemma}{Lemma}[section]
\newtheorem{corollary}{Corollary}[section]
\newtheorem{definition}{Definition}[section]
\newtheorem{proposition}{Proposition}[section]
\newtheorem{theorem}{Theorem}[section]
\newtheorem{notation}{Notation}[section]
\newtheorem{remark}{Remark}[section]
\newtheorem{example}{Example}[section]
\newtheorem{ques}{Question}[section]
\newtheorem{problem}{Problem}[section]
\newtheorem{conjecture}{Conjecture}[section]
\usepackage[left=0.5in,right=0.5in,top=1.5cm,bottom=1.5cm]{geometry}
\usepackage{fancyhdr}
\pagestyle{fancy}
\fancyhf{}
\addtolength{\headheight}{0pt}% obsolete
\lhead{\itshape \scriptsize \chaptername~\thechapter}
\rhead{\itshape \scriptsize \nouppercase{\leftmark}} %\nouppercase !
\renewcommand{\chaptermark}[1]{\markboth{#1}{}}
\cfoot{\thepage}

\title{\href{https://www.romsoc.eu/}{ROMSOC}}
\author{Collector: Nguyen Quan Ba Hong}
\date{\today}

\begin{document}
\maketitle
\setcounter{secnumdepth}{6}
\setcounter{tocdepth}{6}
\tableofcontents

%------------------------------------------------------------------------------%

\chapter{\href{https://www.romsoc.eu/}{Home}}

Product development today is increasingly based on simulation and optimization of virtual products and processes.

Mathematical models serve as digital twins of the real products and processes and are the basis for optimization and control of design and functionality.

The models have to meet very different requirements: Deeply refined mathematical models are required to understand and simulate the true physical processes, while less refined models are the prerequisites to handle the complexity of control and optimization.

To achieve best performance of mathematical modeling, simulation and optimization techniques (MSO), in particular in the industrial environment, it would be ideal to create a complete model hierarchy.

%
The currently most favored way in industrial applications to achieve such a model of hierarchy is to use a sufficiently fine parameterized model and then apply model order reduction (MOR) techniques to tune this fine level to the accuracy, complexity and computational speed needed in simulation and parameter optimization.

%
Although the mathematical models differ strongly in different applications and industrial sectors, there is a common framework via an appropriate representation of the physical model.

%
The main objective of the ROMSOC project is to further develop this common framework and, driven by industrial applications as optical and electronic systems, material engineering, or economic processes, to lift mathematical MSO and MOR to a new level of quality.

The development of high dimensional and coupled systems presents a major challenge for simulation and optimization and requires new MOR techniques.

%
ROMSOC is a \textit{European Industrial Doctorate} (EID) project that will run for 4 years bringing together 15 international academic institutions and 11 industry partners.

It supports the recruitment of 11 Early Stage Researchers (ESRs) working on individual research projects.

In order to train the 11 young researchers for the challenges of multi-disciplinary and international co-operation, their scientific work is embedded in a jointly organized doctoral program, in which lectures and workshops address both scientific content and soft skills.

The researchers are supervised by expert tandems, each consisting of an academic and an industrial representative.

They spend at least half the time in a company, the rest in a research facility.

%
European Industrial Doctorate Programs (EIDs) are sponsored by the European Commission in the framework of Horizon 2020.

They are one type of ,,Innovative Training Networks'', which are part of the Marie Skłodowska-Curie Actions (MSCA) that support the carriers of scientists and encourage their transnational, intersectoral and interdisciplinary mobility.

\begin{remark}
    This project has received funding from the European Union's Horizon 2020 research and innovation programme under the Marie Skłodowska-Curie Grant Agreement No. 765374.
    
    %
    This website reflects the views of the author(s) and does not necessarily reflect the views or policy of the European Commission.
    
    The REA cannot be held responsible for any use that may be made of the information this website contains.
\end{remark}
\href{https://twitter.com/ROMSOC_ITN}{Twitter/ROMSOC ITN}.

\href{https://www.facebook.com/ProjectROMSOC}{Facebook/ProjectROMSOC}.

\section{Imprint}
\begin{itemize}
    \item \textbf{Provider.} The provider of this webpage is the \href{https://www.romsoc.eu/consortium/}{ROMSOC-consortium} consisting of 29 academic and industrial institutions.
    \item \textbf{Speaker \& main contact.} \href{https://www.romsoc.eu/contact/}{Prof. Dr. Volker Mehrmann} at Technische Universität Berlin, Straße des 17. Juni 136, D-10623 Berlin
    \item \textbf{Copyright.} All graphics and contributions used are protected by Copyright.
    
    Any use beyond the implied license to view the site with a web browser requires a prior written consent.
    \item \textbf{Liability.} The content provided on the ROMSOC website have been reviewed with great care.
    
    However, ROMSOC cannot guarantee the accuracy, completeness, or validity of the information provided.
    
    Any liability on the part of ROMSOC is therefore expressly excluded.
    
    Links to external websites have been carefully chosen.
    
    However, since ROMSOC has no control over the contents of external websites, ROMSOC assumes no responsibility for them.
    
    ROMSOC reserves the right to change, amend or delete parts or the entirety of its website content or to suspend its publication temorarily or permanently, both without explicit notice.
    \item \textbf{Privacy.} This website does not use Coockies.
    
    Server Logs (IP, date $+$ time, pages visited) are saved for the duration of 4 weeks and only for the purpose of fault analysis.
    
    After 4 weeks they get deleted automatically and permanently.
    \item \textbf{Responsible Editors.}
    
    Scientific Content: Project Leaders of Project 1--11
    
    Non-scientific Content /Layout /Maintenance:
    
    WDM-Statistics e. U., Mitterschlag 56, 3921 Langschlag, Austria
\end{itemize}

%------------------------------------------------------------------------------%

\chapter{\href{https://www.romsoc.eu/projects/}{Projects}}

\section{Project 1}

\section{Project 2}

\section{Project 3}

\section{Project 4}

\section{Project 5}

\section{Project 6}

\section{Project 7}

\section{Project 8}

\section{Project 9}

\section{Project 10}

\section{\href{https://www.romsoc.eu/optimal-shape-design-of-air-ducts-in-combustion-engines/}{Project 11}}
\begin{center}
    \textsc{Project 11: Optimal Shape Design of Air Ducts in Combustion Engines}
\end{center}

\begin{itemize}
    \item \textbf{Project Partners.}
    \begin{itemize}
        \item Weierstraß Institute for Applied Analysis and Stochastics (WIAS), Germany
        \item Math. Tec GmbH, Austria
    \end{itemize}
    \item \textbf{Project Scientists.}
    \begin{itemize}
        \item Hong Nguyen - Early Stage Researcher (ESR11)
        \item Michael Hintermüller (WIAS, HU Berlin)
        \item Karl Knall (Math. Tec)
    \end{itemize}
    \item \textbf{Project Description.} Many optimal design problems in engineering or biomedical sciences require to determine an optimal shape of a region of interest in order to minimize a number of suitable objectives subject to fluid flow.
    
    Often additional geometric constraints impose further restrictions on the possible design.
    
    %
    Mathematically, this problem results in a constrained multi-objective free form shape optimization problem subject to the Navier-Stokes system.
    
    The goal of this project is the derivation of adjoint based representations of shape  gradient-related descent directions and the numerical realization of associated minimization schemes.
    
    For this purpose and for reasons of efficiency reduced order models (e.g. based on shape-aware adaptive discretization) need to be developed, appropriate primal and dual turbulence models need to be implemented and proper preconditioning of saddle point problems has to be developed.
    
    In this context, applications in the automotive industry and quantitative biomedicine will be addressed.
\end{itemize}

%------------------------------------------------------------------------------%

\chapter{\href{https://www.romsoc.eu/contact/}{Contact}}

\begin{itemize}
    \item \textbf{Project Coordinator.} Prof. Dr. Volker Mehrmann
    
    Technische Universität Berlin
    
    Institut für Mathematik
    
    Straße des 17. Juni 136
    
    D-10623 Berlin
    
    Büro: Raum MA 466
    
    Tel: +49 (0)30 314 - 25736
    
    Fax: +49 (0)30 314 - 79706
    \item \textbf{Project Manager.} Dr. Lena Scholz
    
    Technische Universität Berlin
    
    Institut für Mathematik
    
    Straße des 17. Juni 136
    
    D-10623 Berlin
    
    Büro: Raum MA 463
    
    Tel: +49 (0)30 314 - 79177
    
    Fax: +49 (0)30 314 - 79706
\end{itemize}

%------------------------------------------------------------------------------%

\chapter{\href{https://www.romsoc.eu/205-2/}{Training Activities}}

%------------------------------------------------------------------------------%

\chapter{\href{https://www.romsoc.eu/dissemination/}{Dissemination Activities}}

%------------------------------------------------------------------------------%

\chapter{\href{https://www.romsoc.eu/blog/}{Blog}}

%------------------------------------------------------------------------------%

\chapter{\href{https://www.romsoc.eu/consortium/}{Consortium}}

%------------------------------------------------------------------------------%

\chapter{\href{https://www.romsoc.eu/wim2019/}{WIM2019}}

%------------------------------------------------------------------------------%

\printbibliography[heading=bibintoc]
\end{document}